% vim: set textwidth=78 autoindent:

\section{Weitere Kern-Plugins}

% when the revision of a chapter has been finalized, 
% comment out the following line:
% \updatedisclaimer

Die weiteren Kern-Plugins, die nicht in dem Kapitel\ref{sec:core_plugins},
sondern aus didaktischen Gr�nden an anderer Stelle in diesem Handbuch
beschrieben werden, sind in Tabelle~\ref{tab:other_core} aufgelistet. Dazu
finden Sie jeweils eine Referenz zu dem Kapitel, in dem sie vorkommen.

\begin{table}[H]
\centering
\caption{Weitere Kern-Plugins}\label{tab:other_core}\medskip
\small
 \begin{tabular}{|l|l|p{4in}|}
\hline \textbf{Icon} & \textbf{Plugin} & \textbf{Referenzkapitel}\\
\hline
\includegraphics[width=0.6cm]{diagram_overlay}
 & Diagram Overlay \index{Plugins!Diagram Overlay}& Kapitel~\ref{sec:diagram}\\
\hline
\includegraphics[width=0.6cm]{grass}
 & GRASS \index{Plugins!GRASS Integration} & Kapitel~\ref{sec:grass} und Appendix~\ref{appdx_grass_toolbox_modules}\\
 \hline
\includegraphics[width=0.6cm]{plugin_installer}
 & Plugin Installer \index{Plugins!Plugin Installer} & Kapitel~\ref{sec:python_plugin_installer}\\
\hline
\includegraphics[width=0.6cm]{spiticon}
 & SPIT \index{Plugins!SPIT}& Kapitel~\ref{sec:loading_postgis_data} \\
 \hline
\includegraphics[width=0.6cm]{mIconAddWfsLayer}
 & WFS & Kapitel~\ref{sec:ogc-wfs} \\
\hline
\end{tabular}
\end{table}

